% O008 compact-spectral/SVD bridge -- separately authored component.
% Model provenance: OpenAI Codex gpt-5.6-sol, Ultra, atas arahan pengguna.
% Component license: CC BY-SA 4.0.
% Not authored or endorsed by John M. Erdman or Portland State University.

\chapter{JEMBATAN SPEKTRAL-KOMPAK DAN NILAI SINGULAR}
\label{o008bridge:compact-spectral-svd}

\begin{center}
\small
Komponen asli terpisah, CC BY-SA 4.0. Ditulis dengan bantuan OpenAI Codex
gpt-5.6-sol, Ultra, atas arahan pengguna. Komponen ini bukan tulisan John M.
Erdman dan tidak menyiratkan dukungan beliau atau Portland State University.
\end{center}

Bab~\ref{chap_cpt_ops} telah mengembangkan operator kompak dan dekomposisi
polar, sedangkan Bab~\ref{spectrum} telah mengembangkan spektrum. Bab tentang
teori Fredholm kemudian membuktikan alternatif Fredholm bagi operator
Riesz--Schauder, sifat jangkauan tertutup, serta teori indeks. Tujuan jembatan
ini bukan mengulang teori Fredholm tersebut, melainkan menghubungkan hasilnya
ke geometri eigenvektor operator kompak dan selanjutnya ke dekomposisi nilai
singular.

Semua ruang dalam bab ini merupakan ruang kompleks. Simbol $H_1$ dan $H_2$
menyatakan ruang Hilbert, sedangkan $B$ menyatakan ruang Banach. Jika
$u\in H_2$ dan $v\in H_1$, kita memakai notasi peringkat-satu
\[
  u\otimes v\colon H_1\longrightarrow H_2,
  \qquad
  (u\otimes v)x=\langle x,v\rangle u.
\]

\section{Konsekuensi spektral Riesz--Schauder}

% O008-BRIDGE-ID: O008-BRIDGE-CS-DEF-001
\begin{defn}\label{o008bridge:spectral-multiplicity}
Jika $T\in\ofml B(B)$ dan $\lambda\in\C$, maka
\df{ruang eigen} untuk~$\lambda$ adalah $\ker(T-\lambda I)$. Dimensinya
disebut \df{multiplikitas geometrik} nilai eigen~$\lambda$.
\end{defn}

% O008-BRIDGE-ID: O008-BRIDGE-CS-THM-001
\begin{thm}[Konsekuensi spektral Riesz--Schauder]
\label{o008bridge:riesz-schauder-spectrum}
Misalkan $B$ ruang Banach kompleks berdimensi tak hingga dan
$K\in\ofml K(B)$. Maka:
\begin{enumerate}[label=\rm{(\alph*)}]
  \item setiap $\lambda\in\sigma(K)\setminus\{0\}$ merupakan nilai eigen;
  \item setiap ruang eigen untuk $\lambda\neq0$ berdimensi hingga;
  \item untuk setiap $r>0$, hanya ada berhingga banyak nilai eigen berbeda
        yang memenuhi $|\lambda|\geq r$;
  \item $\sigma(K)\setminus\{0\}$ paling banyak terhitung, setiap titiknya
        terisolasi dalam spektrum, dan satu-satunya titik akumulasi yang
        mungkin adalah~$0$;
  \item $0\in\sigma(K)$.
\end{enumerate}
\end{thm}

\begin{proof}
Ambil $\lambda\neq0$. Operator
$\lambda I-K$ adalah operator Riesz--Schauder dalam pengertian
Bab~15. Alternatif Fredholm~\ref{0040343} menyatakan bahwa operator itu
injektif jika dan hanya jika surjektif. Jika injektif, ia bijektif dan teorema
pemetaan terbuka memberi invers terbatas. Jadi, bila
$\lambda\in\sigma(K)$, operator $\lambda I-K$ tidak injektif; dengan kata
lain, $\lambda$ adalah nilai eigen. Alternatif Fredholm yang sama juga memberi
\[
 \dim\ker(\lambda I-K)<\infty,
\]
yang membuktikan (a) dan~(b) tanpa mengulang argumen Fredholm.

Untuk (c), andaikan terdapat nilai eigen berbeda
$\lambda_1,\lambda_2,\ldots$ dengan $|\lambda_n|\geq r>0$, dan pilih
eigenvektor $x_n\neq0$. Eigenvektor untuk nilai eigen berbeda bebas linear.
Tuliskan $M_n=\spn\{x_1,\ldots,x_n\}$. Menurut lema Riesz, kita dapat memilih
$y_n\in M_n$ dengan $\|y_n\|=1$ dan
\[
  \dist(y_n,M_{n-1})>\frac12.
\]
Karena $(K-\lambda_n I)M_n\subseteq M_{n-1}$, berlaku
$y_n-\lambda_n^{-1}Ky_n\in M_{n-1}$. Untuk $m<n$ juga
$\lambda_m^{-1}Ky_m\in M_m\subseteq M_{n-1}$. Akibatnya
\[
 \left\|\lambda_n^{-1}Ky_n-\lambda_m^{-1}Ky_m\right\|>\frac12.
\]
Namun barisan $(\lambda_n^{-1}y_n)$ terbatas oleh~$1/r$, sehingga
kekompakan~$K$ mengharuskan $(\lambda_n^{-1}Ky_n)$ mempunyai subbarisan
konvergen. Ini bertentangan dengan pemisahan di atas. Jadi (c) benar.

Terakhir,
\[
 \sigma(K)\setminus\{0\}
 =\bigcup_{m=1}^{\infty}
   \{\lambda\in\sigma(K):|\lambda|\geq1/m\}
\]
adalah gabungan terhitung himpunan-himpunan berhingga. Pernyataan isolasi dan
akumulasi mengikuti langsung dari (c), sehingga (d) terbukti. Jika
$0\notin\sigma(K)$, maka $K$ invertibel dan
$I=K^{-1}K$ kompak. Hal ini mustahil pada ruang Banach berdimensi tak hingga:
lema Riesz menghasilkan barisan vektor satuan yang tidak mempunyai
subbarisan konvergen. Jadi operator identitas tidak kompak, dan (e) terbukti.
\end{proof}

% O008-BRIDGE-ID: O008-BRIDGE-CS-REM-001
\begin{rem}\label{o008bridge:no-fredholm-duplication}
Teorema~\ref{o008bridge:riesz-schauder-spectrum} memakai alternatif Fredholm
yang sudah terdapat dalam Bab~15 hanya pada dua titik: kesetaraan
injektif--surjektif dan keberhinggaan dimensi kernel. Jembatan ini tidak
mengulang pembentukan aljabar Calkin, teorema Atkinson, indeks Fredholm, atau
klasifikasi komponen lintasan operator Fredholm.
\end{rem}

% O008-BRIDGE-ID: O008-BRIDGE-CS-EXAM-001
\begin{exam}\label{o008bridge:diagonal-compact-example}
Pada $\ell^2$, tetapkan
\[
 K(x_1,x_2,\ldots)=
   \left(x_1,\frac{x_2}{2},\frac{x_3}{3},\ldots\right).
\]
Operator ini kompak dan swaadjoin. Nilai eigennya adalah $1/n$, masing-masing
bermultiplikitas satu, dan
$\sigma(K)=\{0\}\cup\{1/n:n\in\N\}$. Contoh ini memperlihatkan bahwa nol
memang dapat menjadi titik akumulasi dan bahwa jangkauan operator kompak tidak
harus tertutup.
\end{exam}

\section{Teorema spektral kompak swaadjoin}

% O008-BRIDGE-ID: O008-BRIDGE-CS-LEM-001
\begin{lem}\label{o008bridge:max-eigenvalue}
Jika $T\in\ofml K(H)$ swaadjoin dan $T\neq0$, maka $T$ mempunyai nilai eigen
real $\lambda$ dengan $|\lambda|=\|T\|$.
\end{lem}

\begin{proof}
Untuk operator swaadjoin, hasil tentang jangkauan numerik di Bab~5 memberi
\[
 \|T\|=\sup_{\|x\|=1}|\langle Tx,x\rangle|.
\]
Pilih vektor satuan $x_n$ sehingga nilai mutlak di ruas kanan menuju
$\|T\|$. Karena bentuk kuadratik swaadjoin bernilai real, setelah mengambil
subbarisan terdapat $\alpha\in\{\|T\|,-\|T\|\}$ dengan
$\langle Tx_n,x_n\rangle\to\alpha$. Untuk $\alpha=\|T\|$, misalnya,
\begin{align*}
 \|(T-\alpha I)x_n\|^2
 &=\|Tx_n\|^2-2\alpha\langle Tx_n,x_n\rangle+\alpha^2\\
 &\leq2\alpha\bigl(\alpha-\langle Tx_n,x_n\rangle\bigr)
 \longrightarrow0.
\end{align*}
Kasus $\alpha=-\|T\|$ sama dengan memakai $T-\alpha I=T+\|T\|I$.
Kekompakan~$T$ memberi subbarisan $(Tx_n)$ yang konvergen. Karena
$(T-\alpha I)x_n\to0$ dan $\alpha\neq0$, subbarisan $(x_n)$ sendiri
konvergen ke suatu vektor satuan~$x$. Dengan mengambil limit diperoleh
$Tx=\alpha x$. Jadi $\lambda=\alpha$ adalah nilai eigen real dan
$|\lambda|=\|T\|$.
\end{proof}

% O008-BRIDGE-ID: O008-BRIDGE-CS-LEM-002
\begin{lem}\label{o008bridge:orthogonal-eigenspaces}
Jika $T$ swaadjoin, maka ruang eigen untuk dua nilai eigen berbeda saling
ortogonal. Selain itu, komplemen ortogonal setiap ruang eigen invarian
terhadap~$T$.
\end{lem}

\begin{proof}
Misalkan $Tx=\lambda x$ dan $Ty=\mu y$. Nilai $\lambda$ dan $\mu$ real, dan
\[
 \lambda\langle x,y\rangle
 =\langle Tx,y\rangle
 =\langle x,Ty\rangle
 =\mu\langle x,y\rangle.
\]
Jika $\lambda\neq\mu$, maka $\langle x,y\rangle=0$. Jika $z$ tegak lurus
ruang eigen $E_\lambda$, maka untuk setiap $x\in E_\lambda$,
\[
 \langle Tz,x\rangle=\langle z,Tx\rangle
 =\lambda\langle z,x\rangle=0,
\]
sehingga $Tz\in E_\lambda^\perp$.
\end{proof}

% O008-BRIDGE-ID: O008-BRIDGE-CS-THM-002
\begin{thm}[Teorema spektral untuk operator kompak swaadjoin]
\label{o008bridge:compact-selfadjoint-spectral-theorem}
Misalkan $T\in\ofml K(H)$ swaadjoin. Terdapat himpunan ortonormal
$(e_n)$ yang terdiri atas eigenvektor untuk nilai eigen tak nol
$(\lambda_n)$ sedemikian sehingga
\begin{enumerate}[label=\rm{(\alph*)}]
  \item setiap nilai eigen tak nol diulang menurut multiplikitasnya yang
        hingga;
  \item bila daftarnya tak hingga, $|\lambda_n|\to0$;
  \item $H=\ker T\oplus\clo{\spn\{e_n:n\geq1\}}$;
  \item untuk setiap $x\in H$,
        \[
          Tx=\sum_n\lambda_n\langle x,e_n\rangle e_n;
        \]
  \item jika daftar tak hingga, tetapkan
        $T_N=\sum_{n=1}^N\lambda_n(e_n\otimes e_n)$; jika daftar berhingga
        memiliki panjang~$m$, tetapkan
        $T_N=\sum_{n=1}^{\min\{N,m\}}\lambda_n(e_n\otimes e_n)$.
        Dengan urutan $|\lambda_1|\geq|\lambda_2|\geq\cdots$, berlaku
        $\|T-T_N\|=\sup_{n>N}|\lambda_n|\to0$, dengan supremum himpunan
        indeks kosong didefinisikan sebagai~$0$.
\end{enumerate}
Dengan menambahkan sembarang basis ortonormal bagi $\ker T$, kita memperoleh
basis ortonormal seluruh~$H$ yang terdiri atas eigenvektor~$T$.
\end{thm}

\begin{proof}
Jika $T=0$, ambil daftar kosong dan hasilnya langsung benar. Jika
$\dim H<\infty$, teorema spektral kompleks berdimensi hingga dari
Bab~1, Teorema~\ref{00017}, memberikan basis ortonormal eigenvektor; setelah
eigenvektor bernilai eigen nol dipindahkan ke $\ker T$, semua butir di atas
mengikuti, termasuk konvensi ekor kosong pada~(e). Karena itu, selanjutnya
andaikan $\dim H=\infty$ dan $T\neq0$.
Teorema~\ref{o008bridge:riesz-schauder-spectrum} mengatakan bahwa himpunan
nilai eigennya yang tak nol paling banyak terhitung dan setiap ruang eigennya
berdimensi hingga. Pilih satu basis ortonormal dari setiap ruang eigen tak nol.
Menurut Lemma~\ref{o008bridge:orthogonal-eigenspaces}, gabungan semua basis
tersebut ortonormal, sehingga dapat disusun sebagai barisan $(e_n)$ dengan
nilai eigen terkait~$(\lambda_n)$. Teorema Riesz--Schauder yang sama memberi
$|\lambda_n|\to0$ bila barisan itu tak hingga.

Misalkan $M=\clo{\spn\{e_n:n\geq1\}}$. Komplemen $M^\perp$ invarian. Jika
$T|_{M^\perp}$ tidak nol, Lemma~\ref{o008bridge:max-eigenvalue} akan memberi
eigenvektor tak nol lain yang tegak lurus semua~$e_n$, bertentangan dengan
pemilihan seluruh ruang eigen tak nol. Jadi $T|_{M^\perp}=0$ dan
$M^\perp\subseteq\ker T$. Sebaliknya, setiap vektor dalam $\ker T$ tegak
lurus semua ruang eigen dengan nilai eigen tak nol, sehingga
$\ker T=M^\perp$. Ini membuktikan (c).

Uraikan $x=x_0+\sum_n\langle x,e_n\rangle e_n$ dengan
$x_0\in\ker T$. Kontinuitas $T$ memberi deret pada (d). Untuk ekor operator,
ortonormalitas memberi
\[
 \|(T-T_N)x\|^2
 =\sum_{n>N}|\lambda_n|^2|\langle x,e_n\rangle|^2
 \leq\left(\sup_{n>N}|\lambda_n|^2\right)\|x\|^2.
\]
Jadi $\|T-T_N\|\leq\sup_{n>N}|\lambda_n|$. Arah sebaliknya diperoleh dengan
menerapkan $T-T_N$ pada setiap $e_n$ dengan $n>N$ dan mengambil supremum.
Maka (e) benar dan $T$ merupakan limit norma operator-operator spektral
berperingkat hingga tersebut.
\end{proof}

% O008-BRIDGE-ID: O008-BRIDGE-CS-COR-001
\begin{cor}\label{o008bridge:positive-compact-spectral}
Jika $T$ juga positif dan $T\neq0$, maka semua $\lambda_n>0$ dan
\[
 \|T\|=\lambda_1,
 \qquad
 Tx=\sum_n\lambda_n\langle x,e_n\rangle e_n.
\]
\end{cor}

\begin{proof}
Untuk eigenvektor satuan~$e_n$,
$\lambda_n=\langle Te_n,e_n\rangle\geq0$. Karena daftar hanya memuat nilai
eigen tak nol, semuanya positif. Pernyataan norma mengikuti pengurutan dalam
teorema.
\end{proof}

\section{Nilai singular dan SVD kompak}

% O008-BRIDGE-ID: O008-BRIDGE-CS-DEF-002
\begin{defn}\label{o008bridge:absolute-value-singular-values}
Untuk $T\in\ofml K(H_1,H_2)$, tetapkan
\[
 |T|=(T^*T)^{1/2}.
\]
Nilai eigen positif $s_n$ dari $|T|$, diulang menurut multiplikitas dan
diurutkan tak-menaik, disebut \df{nilai singular}~$T$. Secara ekuivalen,
$s_n^2$ adalah nilai eigen positif dari~$T^*T$.
\end{defn}

Operator $T^*T$ kompak, positif, dan swaadjoin. Karena itu
Teorema~\ref{o008bridge:compact-selfadjoint-spectral-theorem} berlaku pada
$T^*T$. Akar positifnya~$|T|$ juga kompak: melalui kalkulus fungsional,
fungsi akar pada spektrum dapat didekati seragam oleh polinom yang bernilai
nol di~$0$, sehingga $|T|$ merupakan limit norma operator-operator kompak.
Jadi teorema tersebut berlaku pula pada~$|T|$. Kernel ketiga operator terkait oleh
\[
 \ker|T|=\ker(T^*T)=\ker T,
\]
sebab $\langle T^*Tx,x\rangle=\|Tx\|^2$.

Dekomposisi nilai singular untuk operator kompak memperoleh bentuk berikut.

% O008-BRIDGE-ID: O008-BRIDGE-CS-THM-003
\begin{thm}[SVD untuk operator kompak]
\label{o008bridge:compact-svd}
Misalkan $T\in\ofml K(H_1,H_2)$. Terdapat daftar hingga (mungkin kosong) atau
barisan nilai singular $s_1\geq s_2\geq\cdots>0$ dan keluarga ortonormal
$(v_n)$ dalam~$H_1$ serta $(u_n)$ dalam~$H_2$ sedemikian sehingga
\begin{align*}
  T v_n&=s_nu_n, & T^*u_n&=s_nv_n,\\
  Tx&=\sum_n s_n\langle x,v_n\rangle u_n,
      &&x\in H_1.
\end{align*}
Lebih tepat,
\begin{enumerate}[label=\rm{(\alph*)}]
  \item $(v_n)$ merupakan basis ortonormal
        $\clo{\ran T^*}=(\ker T)^\perp$;
  \item $(u_n)$ merupakan basis ortonormal
        $\clo{\ran T}=(\ker T^*)^\perp$;
  \item dalam norma operator,
        \[
          T=\sum_n s_n(u_n\otimes v_n);
        \]
  \item untuk daftar tak hingga tetapkan
        $T_N=\sum_{n=1}^Ns_n(u_n\otimes v_n)$, sedangkan untuk daftar
        berhingga sepanjang~$m$ tetapkan
        $T_N=\sum_{n=1}^{\min\{N,m\}}s_n(u_n\otimes v_n)$; maka
        \[
          \|T-T_N\|=s_{N+1}\longrightarrow0,
        \]
        dengan konvensi $s_{N+1}=0$ setelah daftar berhingga berakhir,
        termasuk untuk daftar kosong ketika $T=0$.
\end{enumerate}
\end{thm}

\begin{proof}
Terapkan teorema spektral kompak swaadjoin pada operator positif~$T^*T$.
Pilih eigenvektor ortonormal $v_n$ dan nilai eigen
$s_n^2>0$ sehingga
\[
 T^*Tv_n=s_n^2v_n.
\]
Tetapkan $u_n=s_n^{-1}Tv_n$. Maka
\[
 \langle u_n,u_m\rangle
 =\frac{1}{s_ns_m}\langle T^*Tv_n,v_m\rangle
 =\delta_{nm},
\]
jadi $(u_n)$ ortonormal, dan
$T^*u_n=s_n^{-1}T^*Tv_n=s_nv_n$.

Teorema spektral memberi bahwa $(v_n)$ merentang
$(\ker(T^*T))^\perp=(\ker T)^\perp$. Karena $T$ nol pada $\ker T$, ekspansi
Fourier setiap $x\in H_1$ menghasilkan
\[
 Tx=\sum_n\langle x,v_n\rangle Tv_n
   =\sum_n s_n\langle x,v_n\rangle u_n.
\]
Jadi $\clo{\ran T}$ termuat dalam rentang tertutup $(u_n)$, sedangkan setiap
$u_n=s_n^{-1}Tv_n$ termasuk dalam $\ran T$; ini membuktikan (a) dan~(b).

Untuk setiap vektor satuan~$x$,
\[
 \|(T-T_N)x\|^2
 =\sum_{n>N}s_n^2|\langle x,v_n\rangle|^2
 \leq s_{N+1}^2.
\]
Kesamaan dicapai pada $x=v_{N+1}$ bila suku itu ada. Maka
$\|T-T_N\|=s_{N+1}\to0$, yang sekaligus membuktikan (c) dan~(d).
\end{proof}

% O008-BRIDGE-ID: O008-BRIDGE-CS-COR-002
\begin{cor}[Pendekatan peringkat hingga terbaik]
\label{o008bridge:best-rank-n}
Jika $R\colon H_1\to H_2$ berperingkat paling banyak~$N$, maka
\[
  \|T-R\|\geq s_{N+1}.
\]
Karena $\|T-T_N\|=s_{N+1}$, pemotongan SVD~$T_N$ merupakan pendekatan
peringkat paling banyak~$N$ yang terbaik dalam norma operator.
\end{cor}

\begin{proof}
Jika daftar nilai singular mempunyai paling banyak~$N$ suku, maka
$s_{N+1}=0$ dan ketaksamaannya langsung. Andaikan $s_{N+1}$ ada.
Subruang $E=\spn\{v_1,\ldots,v_{N+1}\}$ berdimensi~$N+1$, sedangkan
$R|_E$ berperingkat paling banyak~$N$. Jadi terdapat vektor satuan
$x\in E\cap\ker R$. Tuliskan $x=\sum_{j=1}^{N+1}c_jv_j$. Karena
$s_j\geq s_{N+1}$,
\[
 \|(T-R)x\|^2=\|Tx\|^2
 =\sum_{j=1}^{N+1}s_j^2|c_j|^2
 \geq s_{N+1}^2\sum_{j=1}^{N+1}|c_j|^2
 =s_{N+1}^2.
\]
Dengan mengambil supremum atas vektor satuan diperoleh batas bawah yang
diminta; teorema sebelumnya menunjukkan bahwa~$T_N$ mencapainya.
\end{proof}

% O008-BRIDGE-ID: O008-BRIDGE-CS-PROP-001
\begin{prop}[Hubungan dengan dekomposisi polar]
\label{o008bridge:svd-polar}
Dalam dekomposisi polar $T=U|T|$, isometri parsial~$U$ memenuhi
$Uv_n=u_n$ untuk semua indeks singular. Secara khusus,
\[
 |T|=\sum_n s_n(v_n\otimes v_n),
 \qquad
 U|T|=\sum_n s_n(u_n\otimes v_n)=T.
\]
\end{prop}

\begin{proof}
Karena $|T|v_n=s_nv_n$ dan $Tv_n=s_nu_n$, persamaan
$T=U|T|$ memberi $s_nUv_n=s_nu_n$. Nilai $s_n$ positif, sehingga
$Uv_n=u_n$. Kedua ekspansi kemudian mengikuti dari teorema spektral bagi
$|T|$ dan konvergensi norma dalam
Teorema~\ref{o008bridge:compact-svd}.
\end{proof}

% O008-BRIDGE-ID: O008-BRIDGE-CS-EXAM-002
\begin{exam}[Matriks sebagai kasus berhingga]
Jika $H_1=\C^n$ dan $H_2=\C^m$, semua operator otomatis kompak. Setelah
$(v_n)$ dan $(u_n)$ diperluas menjadi basis ortonormal pada domain dan
kodomain, matriks operator mengambil bentuk
\[
  T=U\Sigma V^*,
\]
dengan $U$ dan $V$ uniter serta $\Sigma$ diagonal-persegi-panjang yang entri
tak nolnya adalah $s_1\geq s_2\geq\cdots$. Jadi teorema di atas tepat
memperluas SVD matriks ke operator kompak antar-ruang Hilbert.
\end{exam}

\section{Peta penggunaan}

Urutan konseptualnya bergerak dari operator kompak (Bab~7), melalui spektrum
(Bab~8), menuju Teorema~\ref{o008bridge:riesz-schauder-spectrum}, teorema
spektral kompak swaadjoin, lalu SVD. Prasyarat pembuktiannya lebih rinci:
\begin{enumerate}[label=\rm{(\alph*)}]
  \item pembuktian Riesz--Schauder memakai alternatif Fredholm yang telah
        dibuktikan pada Bab~15;
  \item teorema spektral kompak swaadjoin memakai hasil tersebut bersama
        geometri ruang Hilbert dari Bab~4;
  \item konstruksi $|T|=(T^*T)^{1/2}$ dan SVD juga memakai akar positif dan
        kalkulus fungsional kontinu dari Bab~11.
\end{enumerate}
Peta ini ditempatkan setelah seluruh buku agar tidak menimbulkan ketergantungan
melingkar dan tidak menggandakan pembuktian Fredholm dari Bab~15.
